% !PRINTED 575
\weberpage{575}
But since the form $P$ has none but \gesperrt{rational} coefficients, the latter congruence must also hold to the modulus $p$, and thus
\weq{18}{P(t, t_1+a_1, \ldots, t_n+a_n) \equiv P(t, t_1, \ldots, t_n) \pmod{p}.}
This congruence holds for variable $t, t_1, \ldots, t_n$.

Let us now assume that, contrary to the theorem to be proved, $P(\tau)$ is divisible by $\mathfrak p \mathfrak p_1$; then the congruence
\weq{19}{P(\tau, t_1, \ldots, t_n) \equiv 0 \pmod{\mathfrak p \mathfrak p_1},}
holds, and this remains correct when for the independent variables $t_1, t_2, \ldots, t_n$ the substitution $t_1+a_1, t_2+a_2, \ldots, t_n+a_n$ is made; and since $\tau$ thereby passes into $\tau+\xi$, it follows that
\weq{20}{P(\tau+\xi, t_1+a_1, \ldots, t_n+a_n) \equiv 0 \pmod{\mathfrak p \mathfrak p_1}.}
If, on the other hand, we make the substitution $t=\tau+\xi$ in the congruence (18), it follows that
\weq{21}{P(\tau+\xi, t_1, \ldots, t_n) = P(\tau+\xi) \equiv 0 \pmod{\mathfrak p \mathfrak p_1}.}
Let us now first assume that $\mathfrak p_1$ is different from $\mathfrak p$; then we can reason as follows.

We set $t_1=t_2=\cdots=t_n=0$ in (21), and hence also $\tau=0$. But by this $P(t)$ passes into $t^f$ (apart from multiples of $p$), and from (21) we obtain
\[ \xi^f \equiv 0 \pmod{\mathfrak p \mathfrak p_1}, \]
which, however, contradicts the assumption that $\xi$ is not to be divisible by $\mathfrak p_1$.

But if $\mathfrak p_1 = \mathfrak p$, we arrange (21) according to powers of $\xi$ and obtain
\weq{22}{P(\tau+\xi) = P(\tau) + \xi P'(\tau) + \frac{\xi^2}{2} P''(\tau) \ldots,}
where $P'(t), P''(t), \ldots$ are the derivatives of $P(t)$, and where it is to be observed that the forms
% [sic: Weber prints $P'''(\tau)$ here although the sentence refers to derivatives of $P(t)$.]
\[ \frac{1}{2} P''(t), \ \frac{1}{2.3} P'''(\tau) \ldots \]
are, in spite of the apparent denominators, integral forms in $R$. Now since by (19) and (21) $P(\tau+\xi)$ and $P(\tau)$ are divisible by $\mathfrak p^2$, and likewise, by assumption, $\xi^2, \xi^3, \ldots$, while $\xi$ is not divisible by $\mathfrak p^2$, it follows from (22) that
\[ P'(\tau) \equiv 0 \pmod{\mathfrak p}, \]
from which, on account of
\[ P(t) \equiv \Pi(t) \pmod{\mathfrak p} \]
%CONT>>
